\section{The Mapping: Biology $\leftrightarrow$ Software}\label{sec:mapping}

To compare Agentic Systems and Gene Regulatory Networks (GRNs) under a common typed-interface
abstraction, we map selected components from both domains to a shared mathematical object. We
utilize the category $\mathbf{Poly}$, where objects are polynomial functors representing interfaces, and
morphisms represent interaction protocols. The claim in this section is a correspondence of interface
descriptions, not a proof that the full biological and software domains are equivalent in mechanism,
implementation, or dynamics.

\subsection{Preliminaries: The Category $\mathbf{Poly}$}

In Applied Category Theory, a Polynomial Functor $P$ represents a typed interface for a dynamical system. It is
defined as a sum of representable functors:
\begin{equation}
P(y) = \sum_{o\in O} y^{I_o}.
\label{eq:poly-functor}
\end{equation}
Here, $O$ is the set of possible Positions (or Outputs) the system can expose. For each position $o\in O$, there
is a set $I_o$ of Directions (or Inputs) required to transition the system to a new state.
\begin{itemize}[leftmargin=*]
\item The coefficient $o$ represents the value produced by the system.
\item The exponent $I_o$ represents the capacity to receive information from the environment.
\end{itemize}
This formalism captures the essence of a ``stateful interface'': the system outputs a value $o$ and then waits
for a specific type of input $i\in I_o$ before it can proceed.

\begin{figure}[h]
\centering
\begin{tikzpicture}[>=Latex, node distance=12mm]
  \node[draw, rounded corners, minimum width=30mm, minimum height=10mm] (cap) {Output $o\in O$};
  \node[draw, circle, below left=10mm and 10mm of cap] (i1) {$i_1\in I_o$};
  \node[draw, circle, below=10mm of cap] (i2) {$i_2\in I_o$};
  \node[draw, circle, below right=10mm and 10mm of cap] (i3) {$\cdots$};
  \draw[->] (i1) -- (cap);
  \draw[->] (i2) -- (cap);
  \draw[->] (i3) -- (cap);
  \node[below=12mm of i2] {The Interface $P(y)$};
\end{tikzpicture}
\caption{A visual representation of a Polynomial Functor (often called a ``Mushroom'' or ``Corolla''). The system
offers an Output (the cap) and exposes specific Input ports (the stalks) dependent on that output.}
\end{figure}

\subsection{The Correspondence: Genes and Agent Capabilities}

We now apply this abstract definition to our specific domains.

\begin{definition}[The Gene Object]\label{def:gene-object}
A gene $G$ is a polynomial functor where $O_G$ is the set of expressed proteins and
$I_G=(I_{\text{prot}})_{\text{prot}\in O_G}$ is the \textbf{family} of regulatory-signal sets (transcription factors)
available at each expressed protein:
\begin{equation}
P_{\text{Gene}}(y) = \sum_{\text{prot}\in O_G} y^{I_{\text{prot}}}.
\label{eq:poly-gene}
\end{equation}
\end{definition}

\begin{definition}[The Agent-Capability Object]\label{def:agent-capability}
An agent capability $A$ is a polynomial functor where $O_A$ is the set of generated messages/actions, and
$I_A=(I_{\text{action}})_{\text{action}\in O_A}$ is the \textbf{family} of observation sets available at each action:
\begin{equation}
P_{\text{Agent}}(y) = \sum_{\text{action}\in O_A} y^{I_{\text{action}}}.
\label{eq:poly-agent}
\end{equation}
\end{definition}

\paragraph{Multi-Scale Composition.}
The polynomial functor formalism applies at every level of biological and software organization. Definitions~\ref{def:gene-object} and~\ref{def:agent-capability} establish the \textit{atomic} correspondence: a single gene and a single agent capability share the same interface form. Polynomial functors then compose---via the parallel ($\otimes$), serial ($\circ$), and trace ($\mathrm{Tr}$) operations of \S\ref{sec:operad}---allowing the same interface language to describe progressively richer assemblies:

\begin{center}
\small
\begin{tabular}{@{}lll@{}}
\toprule
\textbf{Biological Scale} & \textbf{Software Scale} & \textbf{Composition}\\
\midrule
Gene & Agent capability & Atomic polynomial functor\\
Cell & Agent runtime & Composite of capabilities $+$ organelles\\
Tissue & Multi-agent subsystem & Wiring diagram of agents (\S\ref{sec:multi-cellular})\\
\bottomrule
\end{tabular}
\end{center}

\noindent A cell is not merely a bag of genes; it is a structured composition with shared infrastructure (organelles) and a boundary (membrane). Likewise, an agent runtime is a structured composition of capabilities with shared components (LLM provider, memory, error handling) and a security boundary. To reduce ambiguity, the rest of the paper uses \textit{capability} for the atomic interface and \textit{agent} for the composed runtime-level object, even though later wiring-diagram examples sometimes use ``agent'' informally for executable boxes. The organelle mappings below describe this cell-level architecture. The multi-cellular organization of \S\ref{sec:multi-cellular} then composes agents into tissues via wiring diagrams---the same formalism, one level up.

\subsection{The Interface: Promoters as Lenses}

In biology, a gene is not universally accessible. It is guarded by a Promoter Region---a specific DNA sequence
that only binds to compatible Transcription Factors. In software, an agent is guarded by an API Schema or
Context Window definition.

We model this gating mechanism using Optics, specifically Lenses. A Lens consists of two maps between a global
state $S$ and a local view $V$:
\begin{enumerate}[leftmargin=*]
\item Get (View): $\mathrm{get}: S \to V$ (Extracting relevant signal from global state).
\item Put (Update): $\mathrm{put}: S\times V \to S$ (Updating global state based on local change).
\end{enumerate}

The ``Promoter'' acts as a filter that determines which part of the global cellular environment ($S$) is visible
($V$) to the gene.
\begin{itemize}[leftmargin=*]
\item \textbf{Biological Lens:} The promoter filters the chaotic cellular soup, allowing the gene to ``see''
only specific molecules (e.g., Lac Repressor).
\item \textbf{Agentic Lens:} The Context Window filters the massive vector database, allowing the agent to
``see'' only the relevant retrieved chunks (RAG).
\end{itemize}

If the input signal does not match the Schema (Promoter), the Lens fails to focus, and the interaction is routed
to an explicit \textbf{inactive/error} case (equivalently, one works with a \textbf{partial} lens, or a total lens
into $V+\mathrm{Error}$) (the agent does not run; the gene is not expressed).

\subsection{Wire-Level Optics: Beyond Lenses}

The Lens formalism models constitutive access: a promoter that either admits or blocks a signal. Biological systems employ richer signal-processing at the interface level. We extend the wiring diagram with two additional optic types from the categorical optics literature.

\paragraph{Prism: Receptor Specificity.}
A membrane receptor does not merely pass or block signals; it selects signals by molecular shape. A Prism on a wire admits values of specific data types and rejects others:
\begin{equation}
\mathrm{prism}_A(\tau, v) =
\begin{cases}
v & \text{if } \tau \in A \\
\bot & \text{if } \tau \notin A
\end{cases}
\label{eq:prism-optic}
\end{equation}
where $A \subseteq T$ is the set of accepted types. This enables fan-out routing: a single output port connects to multiple wires, each guarded by a different prism, directing JSON to one handler and errors to another---analogous to how different receptor types on a cell surface route different ligands to different intracellular pathways.

\paragraph{Traversal: Polymerase Processivity.}
A ribosome does not translate an entire mRNA at once; it processes codons sequentially, applying the same read-translate operation to each element. A Traversal maps a transform $f$ over collection elements on a wire:
\begin{equation}
\mathrm{traversal}_f(\mathbf{x}) = [f(x_i) \mid x_i \in \mathbf{x}]
\label{eq:traversal-optic}
\end{equation}
This models batch processing at the wire level: a list of candidate outputs is transformed element-wise before reaching the downstream agent.

\paragraph{Composition and Coexistence.}
Optics compose: a \textbf{ComposedOptic} applies its constituent optics left-to-right, transmitting only if all accept. A wire may carry both a DenatureFilter (\S5.3) and an Optic, applied in sequence: denaturation strips syntactic structure, then the optic routes or transforms the sanitized content. This layered processing models the biological reality that signal reception involves multiple sequential steps (ligand binding $\to$ receptor conformational change $\to$ intracellular cascade).

\subsection{Epigenetics and State: The Coalgebra}
\label{sec:coalgebra}

Neither genes nor agents are stateless functions. They possess memory.
\begin{itemize}[leftmargin=*]
\item \textbf{Biology:} Epigenetic markers (Methylation, Histone modification) alter how a gene responds to
signals without changing the DNA code.
\item \textbf{Software:} Retrieval Augmented Generation (RAG) and Conversation History alter how an agent
responds to a prompt without changing the LLM weights.
\end{itemize}

We model this as a Coalgebra for the polynomial functor $P$. A dynamical system is defined as a tuple $(S,\phi)$,
where $S$ is the state space and $\phi$ is the structure map:
\begin{equation}
\phi: S \to P(S).
\label{eq:coalgebra-structure}
\end{equation}

By expanding $P(S)$, we derive the two fundamental operations of the state machine:
\begin{enumerate}[leftmargin=*]
\item Readout: $S \to O$ (Given current state/memory, what action do I take?)
\item Update: $\displaystyle \sum_{s\in S} I_{o(s)} \to S$ (Given current state $s$ and a new input
$i\in I_{o(s)}$ compatible with its current output $o(s)$, what is my new state?)
\end{enumerate}

By establishing this formal dictionary (Table~\ref{tab:correspondence-dictionary}), we can compare selected GRN components and agentic
components as instances of the same abstract class of typed dynamical interfaces \textbf{under this
interface-level abstraction}. Proposition~\ref{prop:substrate-independence} below states precisely what this
shared membership does and does not assert.

\paragraph{Composable Coalgebras.}
The coalgebra formalism becomes most useful when made composable. We define two composition operations that mirror the parallel and serial composition of the Operad (\S4):
\begin{itemize}[leftmargin=*]
\item \textbf{Parallel Coalgebra:} Given $(S_1, \alpha_1)$ and $(S_2, \alpha_2)$, their parallel composition has state $S_1 \times S_2$ and applies both readout/update operations to the same input---like two signaling pathways activated by the same ligand.
\item \textbf{Sequential Coalgebra:} Given $(S_1, \alpha_1)$ over $P_1$ and $(S_2, \alpha_2)$ over $P_2$, the sequential composition pipes the readout of the first as input to the second, with joint state $S_1 \times S_2$---modeling signal transduction cascades.
\end{itemize}

\paragraph{Bisimulation: Observational Equivalence.}
For the deterministic state machines considered here, we use the following observational criterion:
two state machines $(S_1, \alpha_1)$ and $(S_2, \alpha_2)$ are equivalent if, for every input sequence,
they produce identical output sequences. We write $M_1 \sim M_2$ when
\begin{equation}
\forall \, \mathbf{i} \in I^*: \quad \mathrm{readout}_1^*(\mathbf{i}) = \mathrm{readout}_2^*(\mathbf{i})
\label{eq:bisimulation}
\end{equation}
where $\mathrm{readout}^*$ denotes the lifted readout over the input sequence. A diverging input (witness) constitutes a proof of non-equivalence. This supports formal comparison of deterministic implementations within the input model considered here; stronger coinductive verification claims are left for future work.

\paragraph{Temporal Coalgebra: Bi-Temporal State.}
The coalgebra above models state at a single point in time. In practice, agents accumulate beliefs that may be corrected retroactively. We extend the state space to carry two independent time indices. Let $\mathcal{T} = (\mathbb{R}_{\geq 0}, \leq)$ denote a totally ordered time domain. A \emph{bi-temporal state} augments the coalgebra with a pair of interval-valued annotations:
\begin{equation}
S_{\mathrm{bt}} \;=\; S \times \underbrace{[\mathcal{T}, \mathcal{T} \cup \{\infty\})}_{\text{valid interval}} \times \underbrace{[\mathcal{T}, \mathcal{T} \cup \{\infty\})}_{\text{record interval}}
\label{eq:bitemporal-state}
\end{equation}
where an open-ended interval $[t, \infty)$ denotes a currently active fact. The key invariant is \emph{append-only correction}: closing a record's transaction interval and appending a new record with a \texttt{supersedes} pointer, rather than mutating the original. This ensures that for any pair $(t_v, t_r)$, the \emph{belief state}---the set of facts valid at $t_v$ and known at $t_r$---is uniquely reconstructible by filtering on both intervals simultaneously.

The readout function becomes time-parameterized: $\mathrm{readout}(s, t_v, t_r)$ returns only those facts whose valid interval contains $t_v$ and whose record interval contains $t_r$. This separates two questions that a single-time coalgebra conflates: ``what is true now?'' (valid-time query) versus ``what did the system believe at decision time?'' (bi-temporal query). The implementation (\S\ref{sec:bitemporal-impl}) instantiates this as \texttt{BiTemporalMemory} with explicit \texttt{retrieve\_valid\_at}, \texttt{retrieve\_known\_at}, and \texttt{retrieve\_belief\_state} methods.

\subsection{The Correspondence, Made Precise}

The dictionary of Table~\ref{tab:correspondence-dictionary} is more than a list of analogies: both genes and
agent capabilities have been presented as objects of a single mathematical kind---polynomial-functor
coalgebras. We name that kind and state exactly what the two domains share.

\begin{definition}[Typed Dynamical Interface]\label{def:typed-interface}
A \emph{typed dynamical interface} is an object of $\mathbf{Coalg}(\mathbf{Poly})$: a pair $(S,\phi)$ with
structure map $\phi\colon S \to P(S)$ for a polynomial functor $P$, equivalently a readout $S \to O$ together
with an update $\sum_{s\in S} I_{o(s)} \to S$ (Eq.~\ref{eq:coalgebra-structure}). A morphism
$(S_1,\phi_1)\to(S_2,\phi_2)$ is a coalgebra homomorphism: a map of state spaces commuting with the structure
maps, hence one that preserves observable behavior in the sense of Eq.~\ref{eq:bisimulation}.
\end{definition}

Definitions~\ref{def:gene-object} and~\ref{def:agent-capability} exhibit the Gene Object and the
Agent-Capability Object as two such interfaces. Following Marom et al.~\cite{marom2026category}---who relate
biological and engineered stimulus--response systems by a structure-preserving functor rather than by
identifying their substrates---we regard the biological and software realizations as two full subcategories of
the same ambient category:
\[
\mathsf{Nat} \subset \mathbf{Coalg}(\mathbf{Poly})\ \ (\text{GRN realizations}),
\qquad
\mathsf{Art} \subset \mathbf{Coalg}(\mathbf{Poly})\ \ (\text{agentic realizations}).
\]

\begin{proposition}[Substrate-Independent Correspondence]\label{prop:substrate-independence}
Each biological and agentic component named in this section admits an interpretation as an object of
$\mathbf{Coalg}(\mathbf{Poly})$. Every row of Table~\ref{tab:correspondence-dictionary} is a pair $(n,a)$ with
$n\in\mathsf{Nat}$ and $a\in\mathsf{Art}$ whose interpretations coincide \emph{as typed dynamical interfaces}:
the same position set $O$ up to relabeling, the same direction family $(I_o)_{o\in O}$, and the same
readout/update shape.
\end{proposition}

\begin{proof}[Proof sketch]
By construction. Definition~\ref{def:gene-object} fixes the $\mathbf{Poly}$ interpretation of each GRN
component and Definition~\ref{def:agent-capability} the agentic one; the rows of
Table~\ref{tab:correspondence-dictionary} are the witnessing pairs, each checked at the level of interface
signature (position set, direction family, and the readout/update maps of
Definition~\ref{def:typed-interface}). The claim is confined to that signature.
\end{proof}

\paragraph{What this asserts, and what it does not.}
The correspondence is a shared \emph{interface} semantics, not an isomorphism. We do not exhibit a mutually
inverse pair of functors between $\mathsf{Nat}$ and $\mathsf{Art}$, and we make no claim that the two domains
agree in mechanism, kinetics, or failure distribution---the empirical sections below in fact document regimes
where the biological structure confers no advantage over a naive baseline. In the words of Marom et al., two
systems that share compositional structure can be related by a functor that ``preserves the interface logic
\ldots without requiring that the two domains share any physical substrate.'' That substrate
independence---not equivalence---is the content of Proposition~\ref{prop:substrate-independence}, and it is
what licenses transporting a design pattern read off a gene-regulatory motif to an agent architecture at all.

\begin{figure}[h]
\centering
\begin{tabular}{@{}c@{\qquad}c@{}}
\begin{tabular}{@{}c@{}}
Gene\\(Function)\\[0.25em]
Transcription Factors ($I$)\\
Proteins ($O$)\\
Promoter Binding\\
Expression
\end{tabular}
&
\begin{tabular}{@{}c@{}}
Agent\\(LLM + Tools)\\[0.25em]
Observations ($I$)\\
Actions ($O$)\\
Schema Match\\
Generation
\end{tabular}
\end{tabular}
\\[0.5em]
\caption{The Structural Correspondence. Both Genes and Agent Capabilities act as transducers converting Input Contexts ($I$) into
Output Expressions ($O$), governed by the same categorical laws (at the level of polynomial-interface models).}
\end{figure}

\begin{table}[h]
\centering
\small
\setlength{\tabcolsep}{4pt}
\begin{tabular}{@{}p{0.27\linewidth}p{0.34\linewidth}p{0.34\linewidth}@{}}
\toprule
Category Concept & Biological Realization (GRN) & Software Realization (Agentic)\\
\midrule
Polynomial Functor ($P$) & Gene Interface & Agent Interface (System Prompt)\\
Output Position ($O$) & Protein Expression & Tool Call / Message\\
Input Direction ($I$) & Transcription Factor Binding & Observation / User Prompt\\
Lens (Optic) & Promoter Region & API Schema / Context Window\\
Internal State ($S$) & Epigenetic Markers (Methylation) & Vector Store / Chat History\\
Morphism ($\circ$) & Signal Transduction Pathway & Data Pipeline\\
\midrule
\multicolumn{3}{@{}l@{}}{\textit{Organelles (Specialized Processing Units)}}\\
\midrule
Template Engine & Ribosome (mRNA $\to$ Protein) & Prompt Template Factory\\
Output Validation & Chaperone (Protein Folding) & Schema Validator / JSON Parser\\
Waste Processing & Lysosome (Autophagy) & Error Handler / Garbage Collector\\
Decision Center & Nucleus (Transcription) & LLM Provider Wrapper\\
Input Filter & Membrane (Immune System) & Prompt Injection Defense\\
Computation Engine & Mitochondria (MIPS) & Runtime Supervisor / Decision Gate\\
\midrule
\multicolumn{3}{@{}l@{}}{\textit{Lifecycle and Rhythms}}\\
\midrule
Lifespan Limit & Telomere Shortening & Operation Counter / Max Iterations\\
Periodic Scheduling & Circadian Oscillator & Health Checks / Heartbeat\\
\bottomrule
\end{tabular}
\caption{The Correspondence Dictionary (Extended). Each row is a witnessing pair $(n,a)$ for
Proposition~\ref{prop:substrate-independence}: a GRN component ($n\in\mathsf{Nat}$) and an agentic component
($a\in\mathsf{Art}$) sharing the same typed-interface signature, not an identity of mechanism.}
\label{tab:correspondence-dictionary}
\end{table}

\begin{table}[h]
\centering
\small
\setlength{\tabcolsep}{4pt}
\begin{tabular}{@{}p{0.27\linewidth}p{0.34\linewidth}p{0.34\linewidth}@{}}
\toprule
\textbf{Biological Concept} & \textbf{Agentic Concept} & \textbf{Formal Structure}\\
\midrule
\multicolumn{3}{@{}l@{}}{\textit{Immunity and Security}}\\
\midrule
Toll-Like Receptor (TLR) & Regex Injection Filter & Pattern Matching\\
MHC Presentation & Provenance Labeling & Functor $\mathcal{P}: \mathbf{Msg} \to \mathbf{Trust}$\\
T-Cell Receptor & Trust Gate & Partial Lens\\
Negative Selection & Injection Training & Penalized Learning\\
Regulatory T-Cell & Confidence Dampening & Suppression Function\\
Immune Memory & Threat Signature Store & Hash-Indexed Cache\\
\midrule
\multicolumn{3}{@{}l@{}}{\textit{Metabolism and Control}}\\
\midrule
ATP & Token Budget & Resource Monoid $\mathcal{R}$\\
mPTP Opening & Fast Apoptosis Trigger & Guard Condition\\
Retrograde Signaling & Phenotype Reshaping & Slow Adaptation\\
Metabolic-Epigenetic State Coupling & Cost-Gated Retrieval Policy & Conditional Access Control\\
\midrule
\multicolumn{3}{@{}l@{}}{\textit{Information and Health}}\\
\midrule
Trophic Factors & Novel Input Signals & Epiplexity $> \delta$\\
Bayesian Brain & Free Energy Minimization & KL Divergence\\
Apoptosis & Agent Termination & State $\to \bot$\\
\midrule
\multicolumn{3}{@{}l@{}}{\textit{Multi-Cellular Organization}}\\
\midrule
Genome & Base Model Weights & Shared Parameters\\
Epigenome & System Prompt + RAG & Phenotype Context\\
Morphogen Gradient & Shared Context Variables & JSON State\\
Morphogen Diffusion & Gradient Propagation on Agent Graph & Discrete PDE on $G$\\
Tissue Boundary & Trust Boundary / Capability Ceiling & Type Barrier\\
Prism Optic & Conditional Wire Routing & Partial Optic\\
Traversal Optic & Batch Wire Transform & Endofunctor on Lists\\
Bisimulation & Agent Equivalence Testing & Observational Equivalence\\
\bottomrule
\end{tabular}
\caption{Extended Correspondence Dictionary: Security, Metabolism, and Organization}
\end{table}

\subsection{Metabolic Coalgebras: Formalizing Resource Constraints}

Finally, we address the physical constraints of computation. Just as biological systems are limited by ATP availability \cite{lynch2015bioenergetic,kempes2017thermodynamic}, agentic systems are limited by token budgets and latency. To model this, we extend our coalgebraic framework to include resource constraints, defining a \textbf{Metabolic Coalgebra}. Mathematically, this is an instance of a Quantitative Coalgebra enriched over a resource monoid, effectively restricting the domain of the state transition function to resource-sufficient states.

We align this definition with the theory of \textbf{Quantitative Polynomial Functors} \cite{nakov2021quantitative}, treating the system as a state machine enriched over a resource monoid.

\begin{definition}[The Resource Monoid]\label{def:resource-monoid}
Let $(\mathcal{R}, +, 0, \ge)$ be an ordered commutative monoid representing computational resources, equipped with a partial subtraction $r \ominus c$ defined whenever $r \ge c$. For a single resource dimension (e.g., token counts), $\mathcal{R} \cong \mathbb{N}$; for multi-dimensional accounting (latency, memory), $\mathcal{R} \cong \mathbb{R}_{\ge 0}^k$.
\end{definition}

\begin{definition}[Metabolic Coalgebra]\label{def:metabolic-coalgebra}
A resource-constrained agent is a coalgebra $(S, \alpha)$ over a polynomial functor $P$, where the state space is the product of the logical state $L$ and the resource state $\mathcal{R}$:
\begin{equation*}
    S \cong L \times \mathcal{R}
\end{equation*}
The structure map $\alpha: S \to P(S) + \bot$ is defined as a \textbf{partial map} guarded by cost. Writing only the continuation state and suppressing the output/readout coordinate of $P(S)$, a transition requiring cost $c \in \mathcal{R}$ has the guarded form:
\begin{equation*}
    \alpha_{\mathrm{cont}}(l, r) =
    \begin{cases}
      (l', r \ominus c) & \text{if } r \ge c \\
      \bot & \text{if } r < c \quad \text{(Apoptosis)}
    \end{cases}
\end{equation*}
\end{definition}

This structure maps to the energetics of transcriptional elongation. A gene or agent capability cannot produce its output instantaneously; it must carry out a sequence of state transitions, each with a cost. The Metabolic Coalgebra models this dependency: if the cellular or computational energy budget is exhausted, execution stalls (Ischemia), and the component fails to execute its function, regardless of its regulatory logic.

This formalism establishes that ``Ischemia'' (Token Starvation) is not merely a runtime error, but a reachable terminal state $\bot$ in the system's dynamics. This mirrors the biological mechanism where a cell that cannot meet its energetic demands undergoes programmed death: energy stress is sensed by pathways such as AMPK, which can in turn engage p53-dependent apoptosis \cite{aubrey2018p53}.


\begin{theorem}[The Metabolic Bound (Qualified)]\label{thm:metabolic-bound}
Assume either a single scalar resource budget $\mathcal{R} \cong \mathbb{N}$ or, more generally, a well-founded scalar potential $\mu: \mathcal{R} \to \mathbb{R}_{\ge 0}$ such that every non-identity transition of cost $c$ satisfies $\mu(r \ominus c) \le \mu(r) - c_{\min}$ for some $c_{\min} > 0$. Assume also that the resource state is monotone nonincreasing (no regeneration), and that infinite executions may not consist solely of stutter / identity steps. Then any execution contains at most $\lfloor \mu(R_{\mathrm{total}}) / c_{\min} \rfloor$ non-identity transitions; in the scalar-budget case this is $\lfloor R_{\mathrm{total}} / c_{\min} \rfloor$. Hence termination is decidable up to stuttering. If regeneration, zero-cost cycles, or unrestricted stutter loops are permitted, termination requires an additional well-founded potential or an explicit budget/termination certificate \cite{boreale2023weighted}.
\end{theorem}

\paragraph{Operational takeaway.}
For an agent-systems reader, the theorem says: if every non-trivial step burns at least a fixed minimum amount of budget and budget never regenerates, then execution length is bounded by initial budget divided by that minimum burn rate. Infinite behavior becomes possible only through free stutter, zero-cost cycles, or regeneration.

\begin{proof}[Proof sketch]
Let the initial resource state be $R_{\mathrm{total}}$. By assumption, every non-identity transition decreases the scalar potential $\mu$ by at least $c_{\min} > 0$. After $N$ non-identity transitions we therefore have
\[
\mu(r_N) \le \mu(R_{\mathrm{total}}) - N c_{\min}.
\]
Because $\mu$ is nonnegative, the right-hand side must remain nonnegative, which implies
\[
N \le \left\lfloor \frac{\mu(R_{\mathrm{total}})}{c_{\min}} \right\rfloor.
\]
In the scalar-budget case this is exactly $\lfloor R_{\mathrm{total}} / c_{\min} \rfloor$. Hence only finitely many non-identity transitions are possible. Under the stated no-regeneration / no-infinite-stutter assumptions, this yields decidability of termination up to stuttering. If those assumptions are removed, the argument no longer forces progress, so an additional well-founded ranking or explicit budget certificate is required.
\end{proof}

\subsection{Additional Organelles: Completing the Cellular Architecture}

Beyond the core interface mapping, biological cells contain specialized organelles that handle distinct aspects of cellular function. We extend the correspondence to four additional structures that map directly to agentic software components.

\paragraph{Ribosome: Template-to-Output Synthesis.}
In biology, the ribosome reads messenger RNA (mRNA) sequences and synthesizes proteins by assembling amino acids according to the genetic code. Transfer RNA (tRNA) molecules carry amino acids to the ribosome, where codons (three-nucleotide sequences) specify which amino acid to add.

In agentic systems, the Ribosome maps to a \textbf{prompt template engine}:
\begin{itemize}[leftmargin=*]
\item \textbf{mRNA} $\to$ Prompt templates with variable slots
\item \textbf{tRNA} $\to$ Context bindings (variable $\to$ value mappings)
\item \textbf{Codons} $\to$ Template directives (variables, conditionals, loops)
\item \textbf{Translation} $\to$ Template rendering with context injection
\end{itemize}
Just as the ribosome ensures that the genetic code is faithfully translated into functional proteins, the software ribosome ensures that abstract prompt templates are instantiated into concrete, well-formed prompts.

\paragraph{Lysosome: Waste Processing and Recycling.}
The lysosome is the cell's recycling center, containing enzymes that break down cellular waste, damaged organelles, and foreign material. Through autophagy, the cell digests its own components to recover building blocks during stress.

In agentic systems, the Lysosome maps to \textbf{error handling and garbage collection}:
\begin{itemize}[leftmargin=*]
\item \textbf{Waste Classification} $\to$ Categorizing failures (timeout, validation error, toxic input)
\item \textbf{Digestion} $\to$ Processing errors to extract debugging information
\item \textbf{Recycling} $\to$ Recovering useful context from failed operations
\item \textbf{Autophagy} $\to$ Periodic cleanup of stale cache and expired state
\item \textbf{Toxic Disposal} $\to$ Secure handling of sensitive data (API keys, PII)
\end{itemize}
The lysosome prevents accumulation of ``cellular debris'' that could poison the system---analogous to memory leaks or error cascades in software.

\paragraph{Nucleus: The Decision Center.}
In eukaryotic cells, the nucleus houses the DNA and serves as the control center for gene expression. Transcription factors enter the nucleus, bind to promoter regions, and initiate transcription of specific genes.

In agentic systems, the Nucleus maps to the \textbf{LLM provider wrapper}:
\begin{itemize}[leftmargin=*]
\item \textbf{DNA} $\to$ Pre-trained model weights (static substrate of capability)
\item \textbf{Transcription} $\to$ Inference (prompt $\to$ response generation)
\item \textbf{Nuclear Envelope} $\to$ Provider abstraction layer (API boundary)
\item \textbf{Nucleolus} $\to$ Tool integration hub (where external capabilities are assembled)
\end{itemize}
The nucleus abstracts the complexity of the underlying LLM, exposing a consistent interface regardless of the provider (Anthropic, OpenAI, Gemini).

\paragraph{Telomere: Lifecycle and Senescence.}
Telomeres are protective caps at the ends of chromosomes that shorten with each cell division. When telomeres become critically short, the cell enters senescence (permanent growth arrest) or apoptosis. The enzyme telomerase can extend telomeres, enabling continued division in stem cells.

In agentic systems, Telomeres map to \textbf{lifecycle management}:
\begin{itemize}[leftmargin=*]
\item \textbf{Telomere Length} $\to$ Remaining operation budget (max iterations)
\item \textbf{Shortening} $\to$ Decrementing counter per operation
\item \textbf{Senescence} $\to$ Graceful degradation (reduced capability mode)
\item \textbf{Apoptosis} $\to$ Clean shutdown when budget exhausted
\item \textbf{Telomerase} $\to$ Renewal mechanism (resetting counters for trusted agents)
\end{itemize}
This provides a biological basis for the common pattern of limiting agent iterations. Rather than arbitrary timeouts, the telomere model frames lifecycle limits as a natural property of the system's ``cellular age.''

\paragraph{Mitochondria: The Metabolic Motherboard.}
Recent scholarship reframes mitochondria not merely as the cell's powerhouse, but as the \textbf{Mitochondrial Information Processing System (MIPS)} \cite{picard2022mips}. Mitochondria sense environmental stress and integrate metabolic signals that shape cellular decision-making. They function as ``social signaling organelles'' that communicate with the nucleus and other organelles.

In agentic systems, the Mitochondria maps to the \textbf{Runtime Supervisor}:
\begin{itemize}[leftmargin=*]
\item \textbf{ATP Production} $\to$ Deterministic computation (tool execution, code evaluation)
\item \textbf{Stress Sensing} $\to$ Monitoring token burn rate vs. informational yield
\item \textbf{Retrograde Signaling} $\to$ Forcing strategy shifts when metabolic efficiency drops
\item \textbf{Fusion/Fission} $\to$ Context fusion under resource constraints
\end{itemize}

\textbf{Temporal Dynamics: Fast vs. Slow Interventions.}
A crucial distinction exists between the timescales of mitochondrial intervention. In biology, retrograde signaling influences nuclear gene expression over hours to days---it is a \textit{developmental} response that reshapes the cell's phenotype. In contrast, acute metabolic stress (ATP depletion, Ca$^{2+}$ overload) triggers immediate responses: cytochrome c release initiates apoptosis within minutes.

We preserve this distinction in the agentic mapping:
\begin{itemize}[leftmargin=*]
\item \textbf{Fast Intervention (Acute):} The Runtime monitors real-time metrics (token velocity, error rate). When thresholds are breached, it triggers immediate Apoptosis---terminating the current chain-of-thought without negotiation. This mirrors the mitochondrial permeability transition pore (mPTP) opening.
\item \textbf{Slow Intervention (Chronic):} The Runtime accumulates statistics across sessions (average efficiency, failure patterns). These inform \textit{Retrograde Responses}---modifications to system prompts, retrieval strategies, or model selection that reshape the agent's ``phenotype'' over deployment cycles. This mirrors how chronic metabolic stress induces mitochondrial biogenesis and metabolic reprogramming.
\end{itemize}
The Runtime does not ``override'' the LLM in the sense of injecting tokens mid-generation; rather, it governs the \textit{boundary conditions} within which generation occurs, and triggers state transitions (continue, pivot, terminate) at defined checkpoints.
